\section*{Capítulo \ref{ch:g12:catalysis} --- Catálisis}

\begin{solution}{exo:g12:catalysis:1}
\omterm{def:g7:chemical-reactions:reactant}{Reactivo}: \ce{H2O2}. Productos: \ce{H2O} y \ce{O2}. \omterm{def:g12:catalysis:catalyst}{Catalizador}: \ce{I-}
(consumido en la etapa 1, recuperado en la etapa 2). Intermedio: \ce{IO-}
(formado en la etapa 1, consumido en la etapa 2).
\end{solution}

\begin{solution}{exo:g12:catalysis:2}
Homogénea; heterogénea; heterogénea; homogénea.
\end{solution}

\begin{solution}{exo:g12:catalysis:3}
El camino sin catalizar. Los niveles inicial y final son los mismos: se
forman los mismos productos, con la misma energía liberada; solo cambia
el camino.
\end{solution}

\begin{solution}{exo:g12:catalysis:4}
Se consume en una etapa y se recupera en otra: aparece en ambos lados de
la suma de las etapas y se simplifica.
\end{solution}

\begin{solution}{exo:g12:catalysis:5}
La reacción ocurre en la superficie: cuanto mayor es la superficie, más
\omterm{def:g7:atoms-and-molecules:molecule}{moléculas} pueden reaccionar a la vez, con la misma masa de \omterm{def:g12:catalysis:catalyst}{catalizador}
(a menudo caro).
\end{solution}

\begin{solution}{exo:g12:catalysis:6}
Suma: \ce{2Fe^{3+} + 2H2O2 + 2Fe^{2+} + 2H+ -> 2Fe^{2+} + O2 + 2H+ +
2Fe^{3+} + 2H2O}. Los \omterm{def:g9:ions:ion}{iones} hierro y los \omterm{def:g9:ions:ion}{iones} hidrógeno aparecen en
ambos lados y se simplifican: \ce{2H2O2 -> 2H2O + O2}.
\end{solution}

\begin{solution}{exo:g12:catalysis:7}
La mitad de \qty{72}{mL} es \qty{36}{mL}: se alcanza al cabo de unos
\qty{35}{min} sin \omterm{def:g12:catalysis:catalyst}{catalizador} y de \qty{3.5}{min} con él. El \omterm{def:g12:catalysis:catalyst}{catalizador}
lo acorta diez veces.
\end{solution}

\begin{solution}{exo:g12:catalysis:8}
La reacción tiene lugar en la superficie del platino; el plomo la cubre y
se queda allí, de modo que los gases ya no pueden llegar al \omterm{def:g9:periodic-table-first-look:metal}{metal}. El
\omterm{def:g12:catalysis:catalyst}{catalizador} está ``envenenado'': presente, pero inútil.
\end{solution}

\begin{solution}{exo:g12:catalysis:9}
Hasta unos \qty{37}{\celsius}, la reacción se acelera con la
temperatura, como cualquier reacción. Por encima, la \omterm{def:g12:catalysis:enzyme}{enzima}, una
proteína, pierde su forma y con ella su actividad; a \qty{80}{\celsius}
queda destruida.
\end{solution}

\begin{solution}{exo:g12:catalysis:10}
\ce{2CO + 2NO -> 2CO2 + N2}: C 2 = 2, O 4 = 4, N 2 = 2. El monóxido de
carbono se oxida con el monóxido de nitrógeno, que se reduce: cada
contaminante elimina al otro.
\end{solution}

\begin{solution}{exo:g12:catalysis:11}
$n(\ce{O2}) = 0.072 / 24.1 = \qty{2.99e-3}{mol}$;
$n(\ce{H2O2}) = 2 \times \num{2.99e-3} = \qty{5.98e-3}{mol}$ en
\qty{10.0}{mL}: $c = \qty{0.598}{mol/L}$.
\end{solution}

\begin{solution}{exo:g12:catalysis:12}
\ce{C2H5OH -> C2H4 + H2O}, una \omterm{def:g12:curly-arrows:categories}{eliminación}; \ce{C2H5OH -> CH3CHO + H2}.
Un \omterm{def:g12:catalysis:selective}{catalizador selectivo} acelera una de varias reacciones posibles mucho
más que las otras, y así decide el producto.
\end{solution}

\begin{solution}{exo:g12:catalysis:13}
El \omterm{def:g12:catalysis:catalyst}{catalizador} no cambia el \omterm{def:g10:reaction-progress-table:system}{estado final}: las curvas con y sin
\omterm{def:g12:catalysis:catalyst}{catalizador} terminan en el mismo volumen de oxígeno. Se obtiene la misma
cantidad de producto, solo que antes.
\end{solution}

\begin{solution}{exo:g12:catalysis:14}
$\num{6.02e23} / \num{5e6} = \qty{1.2e17}{s}$, unos cuatro mil millones
de años. Para hacerlo en un segundo: \num{1.2e17} \omterm{def:g7:atoms-and-molecules:molecule}{moléculas} de \omterm{def:g12:catalysis:enzyme}{enzima}, es
decir, \qty{2e-7}{mol}.
\end{solution}

\begin{solution}{exo:g12:catalysis:15}
Amoníaco \ce{NH3}: $150 \times 17.0 / 14.0 = \qty{182}{Mt}$. \omterm{def:g7:atoms-and-molecules:atom}{Átomos} de
nitrógeno: $\num{1.50e14} / 14.0 = \qty{1.07e13}{mol}$, es decir,
\qty{5.36e12}{mol} de \ce{N2}.
\end{solution}

\begin{solution}{pb:g12:catalysis:1}
\textbf{1.} \ce{2H2O2 -> 2H2O + O2}.

\textbf{2.} Sí: el peróxido de hidrógeno es el \omterm{def:g11:redox:oxidant}{oxidante} del par
\ce{H2O2}/\ce{H2O} y el \omterm{def:g11:redox:oxidant}{reductor} del par \ce{O2}/\ce{H2O2}; se oxida y se
reduce a sí mismo.

\textbf{3.} Sin \omterm{def:g12:catalysis:catalyst}{catalizador}, la reacción es extremadamente lenta.

\textbf{4.} $20 / 22.4 = \qty{0.893}{mol}$.

\textbf{5.} Dos \omterm{def:g10:the-mole:amount}{moles} de peróxido de hidrógeno por \omterm{def:g10:the-mole:amount}{mol} de oxígeno:
\qty{1.79}{mol}.

\textbf{6.} \qty{1.79}{mol/L}.

\textbf{7.} $M = 2 \times 1.0 + 2 \times 16.0 = \qty{34.0}{g/mol}$;
$C_m = 1.79 \times 34.0 = \qty{60.7}{g/L}$.

\textbf{8.} La mitad: \qty{0.893}{mol/L}.

\textbf{9.} La catalasa: una \omterm{def:g12:catalysis:enzyme}{enzima}, disuelta, un \omterm{def:g12:catalysis:catalyst}{catalizador} biológico;
los \omterm{def:g9:ions:ion}{iones} hierro(III): \omterm{def:g12:catalysis:homogeneous}{catálisis homogénea}.

\textbf{10.} \ce{2Fe^{3+} + H2O2 -> 2Fe^{2+} + O2 + 2H+} y
\ce{2Fe^{2+} + H2O2 + 2H+ -> 2Fe^{3+} + 2H2O}; su suma es
\ce{2H2O2 -> 2H2O + O2}.

\textbf{11.} Con \omterm{def:g12:catalysis:catalyst}{catalizador}, el volumen sube mucho más deprisa; las dos
curvas se estabilizan en el mismo volumen final.

\textbf{12.} Los \omterm{def:g9:ions:ion}{iones} hierro(III) (naranjas) se convierten en
hierro(II) (verde pálido) en la primera etapa y vuelven a hierro(III) en
la segunda; al final, todo el hierro vuelve a ser hierro(III).

\textbf{13.} La cocción ha destruido la \omterm{def:g12:catalysis:enzyme}{enzima}.

\textbf{14.} $0.100 \times 20 = \qty{2.0}{L}$.

\textbf{15.} $n(\ce{O2}) = 0.100 \times 0.893 = \qty{0.0893}{mol}$, es
decir, $0.0893 \times 24.1 = \qty{2.15}{L}$.

\textbf{16.} La descomposición lenta libera oxígeno, que acumularía
presión en un frasco hermético.

\textbf{17.} $1.50 / 2 = \qty{0.750}{mol}$ de oxígeno por litro, es
decir, $0.750 \times 22.4 = \qty{16.8}{L}$: ``16.8 volúmenes''.

\textbf{18.} $(1.79 - 1.50)/1.79 \approx \qty{16}{\percent}$.

\textbf{19.} \qty{1.79}{mol/L}.
\end{solution}
